\documentclass[11pt]{article}
\usepackage{amsmath,amssymb,amsthm,hyperref}
\newtheorem{theorem}{Theorem}
\newtheorem{proposition}[theorem]{Proposition}
\newtheorem{lemma}[theorem]{Lemma}
\newtheorem{remark}[theorem]{Remark}
\newcommand{\Q}{\mathbb Q}
\newcommand{\CH}{\operatorname{CH}}
\newcommand{\ch}{\operatorname{ch}}
\newcommand{\Gr}{\operatorname{Gr}}
\newcommand{\End}{\operatorname{End}}
\newcommand{\Hom}{\operatorname{Hom}}
\title{GST selected-block descent, component returns, excess, and integrability}
\author{Mathematics-first continuation; focused self-review}
\date{3 November 2026}
\begin{document}\maketitle

\section{Fixed target and proof contract}
The input is an arbitrary smooth projective complex variety $X$, an integer
$0\leq p\leq\dim X$, and a prescribed
$\alpha\in H^{2p}(X,\Q)\cap H^{p,p}(X)$.
The requested output is a cycle $Z_\alpha\in\CH^p(X)_\Q$ satisfying
$\operatorname{cl}^{\rm can}_X Z_\alpha=\alpha$.
Base checkpoint: \texttt{feb0a99d1e08c4998a01f7239e5bb9e668851354}.
This executes the geometric steps of the whole-space plan, rather than
assuming its provisional Hodge premise. Lean compilation and translation
are outside this work package.

We use the actual affine Jouanolou torsor $q:J\to X$ and the
alpha-dependent holomorphic projector/polynomial approximation already
constructed in the October 23--24 derivations. Work componentwise, so
$J=\operatorname{Spec}R$ is connected, smooth, affine and integral.
Let $K\subset J^{\rm an}$ be a compact finite-cycle carrier such that
$H^{2p}(J,\Q)\to H^{2p}(K,\Q)$ is injective.
Choose it connected, adding finitely many connecting paths when necessary.
All analytic neighborhoods below contain this carrier. A contour gap on
such a neighborhood is an analytic fact, not an algebraic factorization.

\section{An actual full cochain realization of the fold plan}
Here the comparison requested in the plan is constructed for an actual
nerve model of every smooth projective realized world $Y$. This does not
identify an arbitrary preassigned GST graph with that nerve.
Choose a finite good cover $\{U_i\}$ of the compact smooth manifold
$Y^{\rm an}$. Retain its whole nerve, including every nonempty higher
intersection. Let $\mathcal A^b$ be smooth complex-valued $b$-forms and put
\[
 C^{a,b}=\prod_{i_0<\cdots<i_a}\mathcal A^b(U_{i_0\cdots i_a}),
 \qquad D|_{C^{a,b}}=\delta+(-1)^a d.                 \tag{21}
\]
The product omits empty intersections. The Cech differential is the
alternating sum of restrictions. Form degree and nerve degree are distinct.
\begin{theorem}[Filtered nerve--forms comparison and causal defect]
The total complex (21) has $D^2=0$. Restriction of global forms gives a
filtered quasi-isomorphism
\[
 (\mathcal A^\bullet(Y),F^\ell)\longrightarrow
 (\operatorname{Tot}C,F^\ell),\qquad
 F^\ell\mathcal A^b=\bigoplus_{u\geq\ell}\mathcal A^{u,b-u}.       \tag{22}
\]
The rational constant nerve complex $B_\Q^a=\prod_{i_0<\cdots<i_a}\Q$
maps, after complexification, to (21) as constant zero-forms and induces
the canonical comparison with $H^*(Y,\Q)\otimes\mathbb C$.
On cohomology the filtration in (22) is the full Hodge filtration.
Thus every coordinate of the finite-dimensional Betti fiber is retained.
The constant nerve by itself is not endowed with a Hodge filtration by
its vertex labels.

Choose a Kahler metric and let $\Pi$ be harmonic projection, $i$ inclusion
of harmonic forms, and $Q_H$ multiplication by $u-v$ on harmonic $(u,v)$
forms. Then
\[
 \widetilde Q=i Q_H\Pi,\qquad
 d\widetilde Q=\widetilde Qd=0                         \tag{23}
\]
is a genuine cochain lift of the Hodge charge. Its kernel on cohomology
in degree $2p$ is $H^{p,p}$; its kernel on all cochains also contains
acyclic directions and must not be identified with $H^{p,p}$.
For actual degree-preserving cochain transports $T_j$ (or with recorded
equal Tate shifts), the cohomology defect has the exact varying-world
telescoping identity from the plan. In particular for a charge-zero input
$x$, the endpoint is Hodge exactly when that defect annihilates $x$.
\end{theorem}
\begin{proof}
Restrictions commute with exterior differential, so the two mixed terms
in $D^2$ have opposite signs. The other two terms vanish by $\delta^2=0$
and $d^2=0$. For a partition of unity $\{\rho_i\}$ subordinate to the
cover, the augmented Cech complex of forms is contracted in positive
nerve degree by
$(hc)_{i_0\cdots i_{a-1}}=\sum_j\rho_j c_{j i_0\cdots i_{a-1}}$,
with alternating extension in the indices. Products are extended by zero
where necessary. The identity $\delta h+h\delta=1$ in positive degree
follows by expanding and using $\sum\rho_j=1$; at degree zero the kernel
is precisely the restrictions of global forms. Multiplication by $\rho_j$
preserves each form type. Apply this row acyclicity to the bounded
double complex, and separately to every $F^\ell$, to prove (22).
The contraction is a Cech contraction; it is not falsely asserted to
commute with $d$.

Every nonempty good-cover intersection is contractible. The ordinary
Poincare lemma on it makes its de Rham complex a resolution of the
constant scalars. The other bounded spectral sequence consequently proves
the constant-nerve quasi-isomorphism. This is the de Rham/nerve comparison
with rational Betti structure, rather than a declaration of a cycle map.
Compact Kahler Hodge theory identifies $F^\ell$ with the sum of harmonic
types whose first index is at least $\ell$; harmonic representatives give
both inclusions. This foundational identification is the one in Deligne's
official problem description, Section 1, p.1.

Harmonic forms are closed and orthogonal to exact forms, proving (23).
On their degree-$2p$ subspace, $u+v=2p$ and charge $u-v=0$ give $u=v=p$.
For $D_j=Q_jT_j-T_jQ_{j-1}$, expansion of
\[
 \sum_{j=1}^k T_k\cdots T_{j+1}D_jT_{j-1}\cdots T_1
\]
cancels adjacent internal $Q_j$ terms, leaving
$Q_kT_k\cdots T_1-T_k\cdots T_1Q_0$.
All operators here act on the actual cohomology just compared.
The empty path gives zero. Safety at every world requires this test on
each prefix; a final cancellation does not certify the prefixes.
\end{proof}

For two actual smoothly homotopic maps $f_0,f_1:Y\to Z$, the concrete
operator $H\omega=\int_0^1\iota_{\partial t}F^*\omega\,dt$ for their
homotopy $F$ satisfies $f_1^*-f_0^*=dH+Hd$ by Stokes. Thus the required
two-path relation is provided when that geometric homotopy is provided;
reachability alone still supplies none. Higher geometric cells require
their actual parameterized maps and compatible boundary restrictions.
The realization (21) supplies all nerve faces automatically, without
pretending that every unrelated GST relation has such a geometric cell.

Finally a product index $I=\prod_i\operatorname{Fin}(n_i)$ can coordinate
this cohomology only after an actual basis comparison of the same dimension
has been chosen. Then $c\mapsto\sum_{i\in I}c_i\beta_i$ is an isomorphism
and preserves the full amplitude function. The single separated-tensor
locus is not substituted for it. Empty products have one vacuum coordinate;
a zero factor makes the product index empty. Neither a hole count nor a
nerve degree has been identified with $\dim Y$.

\section{A single pure-degree topological encoding}
\begin{lemma}[Encoding without a pair of unrelated cluster problems]
For $p>0$ there is a genuine finite-rank topological bundle $F$ on
$J^{\rm an}$, an integer $\mu>0$, and a topological complement $G$ such that
\[
 F\oplus G\simeq\underline{\mathbb C}^{N},\qquad
 \ch(F)=r+\mu q^*\alpha,\qquad
 \ch(G)=m-\mu q^*\alpha,\quad r+m=N.                    \tag{1}
\]
These are equalities of the full rational even Chern character, with
$\alpha$ placed in degree $2p$. They also hold for non-Hodge even classes.
One may take $r,m>0$. The previous Stein realization and polynomial
approximation therefore give a single regular matrix $a\in M_N(R)$
with a separated rank-$r$ local cluster $F_a$ on a neighborhood $V$ of $K$,
and $\ch_p(F_a)|_K=\mu q^*\alpha|_K$.
\end{lemma}
\begin{proof}
Rational topological Chern character on the finite CW homotopy type of
$J^{\rm an}\simeq X^{\rm an}$ is an isomorphism from rational $K^0$ to
rational even cohomology. Clear denominators for the class having only
the required positive-degree component. Represent the resulting integral
$K^0$ class as $[A]-[B]$. A finite-rank complement $B'$ with
$B\oplus B'$ trivial converts this difference to $[A\oplus B']$ minus
a trivial bundle. Thus the genuine bundle $F=A\oplus B'$ has the required
positive-degree character. Take a finite-rank complement $G$ for $F$.
Adding trivial summands makes the two ranks positive.
Pull these bundles to the Stein realization and use the earlier holomorphic
direct-summand construction. Approximation on $K$ preserves its contour
bundle by the gapped homotopy. This last step is precisely the previously
proved analytic construction; it is not an algebraization theorem.
\end{proof}
If $\alpha\ne0$, (1) forces $r,m\geq p$: Newton's identities give
$c_j(F)=0$ for $0<j<p$ and
$c_p(F)=(-1)^{p-1}(p-1)!\mu q^*\alpha\ne0$, so $r\geq p$;
the same argument applies to $G$. This rank observation will matter for
the unpolarized excess construction.

\section{Global factor descent gives the actual original-X cycle}
Let $\chi_a(t)=\det(tI-a)$. Suppose its discriminant is not identically
zero. On the regular spectral open, a chosen set of $r$ roots is called
globally stable if continuation around every loop preserves that set.
Its matching to $F_a$ is part of the assertion, not an arbitrary choice
of some unrelated factor of the same degree.

\begin{theorem}[Selected factor, coherent kernel, and canonical return]
Suppose the local rank-$r$ cluster of $a$ is represented by a monic
factorization
\[
 \chi_a=g h\quad\hbox{in }\operatorname{Frac}(R)[t],
 \qquad \deg g=r,                                      \tag{2}
\]
whose roots on $V$ are exactly that cluster, with no root shared with $h$
anywhere on $V$. Then $g,h\in R[t]$. The actual coherent sheaf
\[
 \mathcal M_g=\ker\bigl(g(a):\mathcal O_J^N\to\mathcal O_J^N\bigr)
                                                               \tag{3}
\]
is perfect on $J$, is the selected vector bundle on
$J\setminus V(\operatorname{Res}(g,h))$, and the actual rational cycle
\[
 Z_\alpha=\frac1\mu(q^*)^{-1}\ch_p^{\CH}([\mathcal M_g])
       \in\CH^p(X)_\Q                                  \tag{4}
\]
has canonical Betti class $\alpha$. Resultant zeros away from $K$ are
allowed. No vanishing of their boundary classes is required.
\end{theorem}
\begin{proof}
Every root of the monic $\chi_a$ is integral over $R$. The coefficients
of $g,h$, being elementary symmetric functions of their selected roots,
are integral over $R$, and are in $\operatorname{Frac}(R)$ by (2).
Normality of $R$ puts them in $R$. The resultant is not identically zero
since the factors are coprime on $V$. After inverting it, Bezout gives
$u g+v h=1$. Cayley--Hamilton gives $g(a)h(a)=0$. Thus
$e=v(a)h(a)$ is an idempotent with image $\ker g(a)$; indeed
$e(1-e)=v(a)h(a)u(a)g(a)=0$, and $e$ acts as identity on that kernel.
On $V$ it is the Riesz cluster projector, even at collisions internal
to $g$ or to $h$. In particular (3) restricts to $F_a$ near $K$.

The kernel is a finite $R$-module. A smooth finite-dimensional affine
coordinate ring is regular of finite global dimension, so this module
has a finite resolution by finite projective modules and defines a genuine
perfect-complex $K_0$ class. Its Chern character in Chow exists even where
the sheaf is not locally free. The canonical comparison commutes with
restriction and with the vector-bundle Chern character of that resolution.
Consequently its degree-$2p$ Betti character restricts on $K$ to
$\mu q^*\alpha$. Injectivity on $K$ proves equality on all of $J$.
The actual affine-bundle Chow isomorphism and its Betti comparison now
give (4). This argument retains all boundary terms through the coherent
kernel instead of discarding them by localization.
\end{proof}

The globally stable root-set condition is equivalent to (2) for a
generically separable polynomial: in its splitting field the coefficients
of $\prod_{i\in I}(t-\lambda_i)$ are fixed precisely when the root subset
is fixed by the Galois/monodromy action. In the complex finite etale
description, this says that the relevant subset selector has a degree-one
component over the regular region. It does not say every connected
component has degree one. One must match the analytic cluster along
continuation, not just at a single point of a disconnected intersection.

This is a proved construction from an actual matched factor. Neither
Hodge type nor small polynomial idempotency defect has supplied (2).

\section{Compute the individual excess blocks, not only their total}
First let $F\oplus G=\mathcal O_J^N$ be an actual algebraic splitting,
of ranks $r,m>0$, and let $P$ be its projector onto $F$.
For a selector rank $s$, put
\[
 k_{\min}=\max(0,s-m),\quad k_{\max}=\min(r,s),\quad l=s-k.
\]
On $M=\Gr_J(s,F\oplus G)=J\times\Gr(s,N)$ consider the invariant-subspace
section $S\to (F\oplus G)/S$ induced by $P$. Its zero scheme is the disjoint
union
\[
 B_k=\Gr_J(k,F)\times_J\Gr_J(l,G),\qquad
 d_k=k(r-k)+l(m-l).                                    \tag{5}
\]
This is an equality of schemes: over any test algebra, $P$ preserving a
subbundle splits it by the idempotent; its locally constant two ranks
give exactly (5).

\begin{theorem}[Exact excess contribution of a block]
The contribution of $B_k$ to the localized zero-section intersection is
\[
 c_{d_k}(T_{B_k/J})\cap[B_k].                           \tag{6}
\]
For $p>0$, proper push of its product with $\ch_p(S)$ is
\[
 A_k\ch_p(F)+B_k'\ch_p(G),                             \tag{7}
\]
where, with binomial coefficients zero outside their usual range,
\[
 A_k=\binom{r-1}{k-1}\binom m{s-k},\qquad
 B_k'=\binom rk\binom{m-1}{s-k-1}.                     \tag{8}
\]
Since the ambient bundle is trivial, (7) equals
$c_k\ch_p(F)$, where $c_k=A_k-B_k'$. These are integral coefficients,
not unspecified nonzero pairings. They satisfy
\[
 \sum_{k=k_{\min}}^{k_{\max}}c_k=0.                    \tag{9}
\]
When $s=r$, $c_r=1$. Thus the distinguished block returns exactly
$\ch_p(F)$; its cancellation has explicit contributions from the other
blocks. Positive-dimensional blocks use (6), not their unweighted
fundamental cycles.
\end{theorem}
\begin{proof}
On $B_k$, write $S=S_F\oplus S_G$ and the quotient as $Q_F\oplus Q_G$.
The vertical tangent of $M$ splits as the sum of the within-block tangent
$\Hom(S_F,Q_F)\oplus\Hom(S_G,Q_G)=T_{B_k/J}$ and the cross-block normal
$\Hom(S_F,Q_G)\oplus\Hom(S_G,Q_F)$. The differential of the section is
zero on the former and multiplication by $-1,+1$ on the latter.
The normal cone therefore embeds as that cross-block subbundle in the
section bundle, with quotient $T_{B_k/J}$. The excess formula gives (6).
There is no real-orientation minus sign: both normal scalar maps are
complex-linear isomorphisms.

The relative Grassmann Euler-transfer identity is
\[
 (\Gr(k,F)\to J)_*\bigl(\ch_p(S_F)c_{k(r-k)}(T_{\Gr(k,F)/J})\bigr)
   =\binom{r-1}{k-1}\ch_p(F),                          \tag{10}
\]
and the same push without $\ch_p$ is $\binom rk$.
To verify it after splitting $F$ into line bundles, use a constant diagonal
endomorphism with distinct eigenvalues. Its vertical tangent section has
exactly the coordinate-subspace sections as regular zeros. On their union,
$\ch_p(S_F)$ is the sum of the chosen line characters. Each line occurs
in $\binom{r-1}{k-1}$ of them. The splitting principle makes this a universal
Chow identity: on a flag bundle one has a line filtration; deformation
of each extension to its graded sum preserves the identity by homotopy
invariance. The pullback to the flag bundle is injective. This also proves
the end cases $k=0,r$. Apply (10) to the two factors, using
$\ch_p(S)=\ch_p(S_F)+\ch_p(S_G)$, to get (7).
Both sums of (8) are $\binom{N-1}{s-1}$ by counting $s$-subsets containing
one fixed label; subtract them to get (9). The formula for $c_r$ is immediate.
\end{proof}

\section{A finite geometric test on the actual alpha-dependent matrix}
Now return to a regular approximation $a$ and its analytic cluster on $V$.
Take the small generic constant perturbation from the prior regular-incidence
theorem, so the actual invariant-subspace scheme $Z_a\subset J\times\Gr(s,N)$
is smooth of dimension $\dim J$, all its components dominate $J$, and
the section is regular. Put $f:Z_a\to J$.
On $V$ the actual analytic contour projector splits the trivial bundle as
$F_a\oplus G_a$. Every invariant subspace is preserved by this projector,
so $f^{-1}V$ is the disjoint union of open-and-closed analytic pieces $Z_k$
with the two dimensions $k,s-k$.

\begin{lemma}[Local block return for the regular approximation]
In rational cohomology of $V$,
\[
 (f|_{Z_k})_*\ch_p(S)=c_k\ch_p(F_a).                  \tag{11}
\]
These are Betti equalities. No Chow class of the possibly nonalgebraic
analytic bundle $F_a$ is being asserted.
\end{lemma}
\begin{proof}
Use the relative Grassmann product in (5) for the analytic splitting.
Cross-block differentiation of invariance is the Sylvester map
$\phi\mapsto a_G\phi-\phi a_F$ or its reverse. It is invertible because
the two spectra are separated. The remaining section is the product of
the internal invariant-subspace tangent sections on the two relative
Grassmannians. Localized complex Euler classes therefore reduce its proper
push to $\ch_p(S)c_{d_k}(T_{B_k/V})$, regardless of internal collisions.
The regularity of $Z_a$ identifies this localized class with its actual
fundamental class. The cohomological version of (10) gives (11).
This also follows by homotoping the block endomorphisms to the two scalar
blocks; the compact Grassmann fibers ensure properness throughout.
\end{proof}

There is a finite, actual construction of the block-complete packages.
Make a bipartite graph with (a) irreducible components of $Z_a$ and (b)
the integers $k_{\min},\ldots,k_{\max}$; connect a component to $k$ when
its analytic restriction meets $Z_k$. Components of this finite graph
give a partition into packages. For a package let $A$ be its set of $k$'s
and $C_A$ the union of its global algebraic components. Every global
component dominates $J$ and hence meets $V$; every $Z_k$ is nonempty.
By construction
$C_A^{\rm an}\cap f^{-1}V=\coprod_{k\in A}Z_k$.

\begin{theorem}[Computed component-package source]
Set $c_A=\sum_{k\in A}c_k$. The genuine rational Chow class
\[
 W_A=f_*\bigl(\ch_p^{\CH}(S|_{C_A})\cap[C_A]\bigr)
       \in\CH^p(J)_\Q                                \tag{12}
\]
has canonical Betti class $c_A\mu q^*\alpha$. If some package has
$c_A\ne0$, then
\[
 Z_\alpha=(\mu c_A)^{-1}(q^*)^{-1}W_A                 \tag{13}
\]
is the required cycle on the original $X$.
\end{theorem}
\begin{proof}
The components are smooth and disjoint; their union has a genuine Chow
fundamental class and the restricted tautological bundle is algebraic.
Canonical proper functoriality, (11), and addition over its local pieces
compute the restriction of its Betti class to $K$ as
$c_A\ch_p(F_a)|_K=c_A\mu q^*\alpha|_K$.
Injectivity on $K$ and the affine-bundle comparison prove (12)--(13).
\end{proof}

This computes the class of a specific source, rather than assuming a
coupling for an unspecified cycle. It does not prove that a nonzero package
exists. The total package has coefficient zero by (9). Components that cut
individual $Z_k$ into smaller pieces can have finer classes; the package
test neither evaluates nor rules out those finer sources.
The operation in (11)--(12) is proper pushforward of a Chern-character
\emph{Chow cycle}, not the Chern character of a derived sheaf pushforward.
It therefore uses canonical Gysin compatibility directly; no Todd term
from a different $K_0$ pushforward formula has been omitted.

\begin{proposition}[Root-transitive package cancellation]
If the actual regular spectral root monodromy is transitive on the $N$
individual roots, every block-complete package above has $c_A=0$.
This requires neither full symmetric monodromy nor transitivity on
$s$-subsets. If local within-cluster monodromy contains
$S_r\times S_m$, each global component is itself block-complete and
therefore has zero degree-$2p$ canonical return on $J$.
\end{proposition}
\begin{proof}
A package's subsets form a union of actual global monodromy orbits, and
locally a union of the shells $|I\cap\{1,\ldots,r\}|=k$ for $k\in A$.
The number of its subsets containing a fixed label in the first block is
$\sum_{k\in A}A_k$; for a fixed label in the second it is
$\sum_{k\in A}B_k'$. Root transitivity makes these numbers equal, proving
$c_A=0$. If local within-block permutations are present, they are
transitive on each shell, so no $Z_k$ can be split among distinct global
components; then each component is already such a package. No assertion
of these monodromy hypotheses for every approximation is made.
\end{proof}

\section{Execute a nonzero-Hodge-input descent test}
The following fixture has an explicit arbitrary-$X$ implication: it refutes
the proposed universal inference from Hodge type and a target-detecting
contour gap to a globally selected algebraic root block. It is not a
substitute for the universal target.

Let $J_0=\{x^2+yz=1\}\subset\mathbb A^3$ and
\[
 P=\frac12\begin{pmatrix}1+x&y\\z&1-x\end{pmatrix},\qquad
 a_\epsilon=P+\epsilon\begin{pmatrix}1&0\\0&-1\end{pmatrix},
 \quad0<\epsilon<\tfrac12.                            \tag{14}
\]
Here $P^2=P$ and $\operatorname{tr}P=1$. Sending a projector to its image
line identifies $q_0:J_0\to\mathbb P^1$ as the affine-line torsor of
complements to the tautological line. Thus $\operatorname{im}P=q_0^*\mathcal
O(-1)$ has nonzero Chern class, and $J_0$ is the genuine Jouanolou model.
The Hermitian-projector sphere
$K_0=\{x\in\mathbb R,z=\overline y,x^2+|y|^2=1\}$ is a deformation
retract and detects that class.

\begin{proposition}[A genuine algebraic Hodge class with no selected sheet]
For every $\epsilon$ as in (14), the eigencluster near $1$ on a neighborhood
of $K_0$ has the nonzero class $-q_0^*h$. Nevertheless its characteristic
polynomial has no proper factor over $\mathbb C(J_0)$, its global
rank-one incidence is a connected smooth double cover, and the complete
tautological first-character return is zero in $\CH^1(J_0)_\Q$.
\end{proposition}
\begin{proof}
Direct calculation using $x^2+yz=1$ gives
\[
 \chi_{a_\epsilon}(t)=t^2-t-\epsilon x-\epsilon^2,
 \quad D_\epsilon=1+4\epsilon x+4\epsilon^2.           \tag{15}
\]
On $K_0$, $D_\epsilon\geq(1-2\epsilon)^2>0$.
The Hermitian eigenvalues remain separated throughout the homotopy
$a_{t\epsilon}$, $0\leq t\leq1$. Their upper line therefore has the
same nonzero first class as $\operatorname{im}P$.
But $D_\epsilon$ has a simple zero along
\[
 x=-\frac{1+4\epsilon^2}{4\epsilon},\qquad yz=1-x^2\ne0.
\]
This is a smooth irreducible divisor isomorphic to $\mathbb C^*$.
Its valuation in $D_\epsilon$ is one, so $D_\epsilon$ is not a square
in the function field. This proves nonsplitting and connected monodromy.
The double cover $w^2=D_\epsilon$ is smooth since the branch divisor is.
The matrix is never scalar: scalarity would force $y=z=0$, $x=-2\epsilon$,
and hence $4\epsilon^2=1$, excluded. Thus its eigenkernel has constant
dimension one on this cover, identifying it with the full incidence.
The regular-incidence top-Chern calculation, now including the branch
fibers, gives the zero first-character return. Equivalently (8) gives
coefficients $-1,+1$, merged into one global package.
For $\epsilon=1/10$, $D=(26+10x)/25$ and the branch is $x=-13/5$,
strictly outside the detecting sphere. A branch outside that sphere still
obstructs global algebraic selection.
\end{proof}
The actual input here is $X=\mathbb P^1$, $p=1$, $\alpha=-h$, already
algebraic. The result refutes the claimed selection implication, not Hodge
algebraicity. For any original $X$, taking its product with this example
retains the same obstruction on the $\mathbb P^1$ factor.

\section{Execute the correlated-matrix excess route}
Let $B=\Gr(s,N)$, $D=s(N-s)$, and let $a_1,\ldots,a_t$ be actual regular
matrices on $J$, with $t\geq2$. Simultaneous invariance is the actual zero
scheme of a section of $(T_B)^{\oplus t}$ on $J\times B$.
For any nonnegative-degree tautological insertion its \emph{complete
virtual} proper Chow return is zero: the ambient push is the integral of
that insertion times $c_D(T_B)^t$, a class of degree at least $tD>D$
on $B$ when $0<s<N$. This statement is about the complete localized
intersection; an excess component's ordinary fundamental cycle is not
automatically that localized class.

For $a_1=\cdots=a_t=P$ and $s=r$, the distinguished actual component is
the section $J$ corresponding to $F$. Its normal in $J\times B$ is
$E=\Hom(F,G)$ of rank $rm$. Differentiating the $t$ equations gives
a diagonal inclusion $E\hookrightarrow E^{\oplus t}$, up to invertible
signs, so its exact excess contribution is
\[
 c_{(t-1)rm}(E^{\oplus(t-1)})\cap[J].                 \tag{16}
\]
For nonzero pure-degree-$p$ encoding (1), $rm\geq p^2$. Consequently
the unpolarized contribution (16) has codimension greater than $p$ for
$p\geq2$. It cannot itself be the missing codimension-$p$ source.

This dimension obstruction does not rule out dimensional folds. They can
be computed explicitly, rather than silently ignored. For one extra
equation, a projective auxiliary variable and an $\mathcal O(1)$ twist
can extract the coefficient
\[
 \pi_*c_{rm}(E\boxtimes\mathcal O_{\mathbb P^{rm-p}}(1))
   =c_p(E),                                           \tag{17}
\]
where $\pi:J\times\mathbb P^{rm-p}\to J$.
This follows by expanding the top Chern class as
$\sum_{j=0}^{rm}c_j(E)h^{rm-j}$ and integrating $h^{rm-p}$.
Under (1), its canonical degree-$2p$ character is
\[
 \operatorname{cl}^{\rm can}c_p(E)
 =(-1)^{p-1}(p-1)!\,[m(-1)^p-r],\mu q^*\alpha.       \tag{18}
\]
Indeed $\ch(E)=\ch(F^*)\ch(G)$ has positive terms zero below $p$,
and degree $p$ equal to $[m(-1)^p-r]\mu q^*\alpha$; Newton gives (18).
The coefficient is nonzero for odd $p$; for even $p$ stabilization can
make $m\ne r$. Formula (17) is a valid polarized Chern conversion whenever
$F,G$ are actual algebraic bundles. It is not an algebraic construction
from the merely analytic $F_a,G_a$. Moreover an individual section of
$\mathcal O(1)$ vanishes on an auxiliary hyperplane, so any proposed
global matrix realization of this twist must retain its additional
boundary components. Equations (17)--(18) alone do not identify that
realization's selected localized component. The global selection question
has not disappeared in the dimension fold.

\section{Execute the integrable-source and source-deformation test}
The following test applies to every original $X$ via the actual smooth
projective world $Y=X\times A$, where $A=E_1\times E_2$ is a product of
elliptic curves. Pull back two independent constant antiholomorphic
one-forms $\eta_1,\eta_2$ from $A$, and set on the trivial rank-two bundle
\[
 U=\begin{pmatrix}0&1\\0&0\end{pmatrix},\quad
 V=\begin{pmatrix}0&0\\1&0\end{pmatrix},\quad
 \theta=U\eta_1+V\eta_2.                              \tag{19}
\]
\begin{proposition}[An explicit second-order obstruction invisible to Chern type]
The first-order operator $\bar\partial+t\theta$ defines an actual
holomorphic deformation over $\mathbb C[t]/(t^2)$. There is no global
$(0,1)$-form $\psi$ making $\bar\partial+t\theta+t^2\psi$ integrable
over $\mathbb C[t]/(t^3)$, although its underlying smooth bundle is trivial
and all its topological Chern classes are zero, hence of Hodge type.
\end{proposition}
\begin{proof}
Expand the square. The coefficient of $t$ is $\bar\partial\theta=0$.
The coefficient of $t^2$ is
\[
 \bar\partial\psi+\theta\wedge\theta,
 \qquad \theta\wedge\theta=[U,V]\eta_1\wedge\eta_2,
 \qquad [U,V]=\begin{pmatrix}1&0\\0&-1\end{pmatrix}.  \tag{20}
\]
The Dolbeault class $[\eta_1\wedge\eta_2]$ is nonzero on $A$:
wedge with a nonzero constant holomorphic two-form and integrate over
$A$; the integral is nonzero, whereas an exact form would give zero by
Stokes. It remains nonzero after pullback to $Y$, since restriction to
a fiber $\{x\}\times A$ is a left inverse. Its product with the nonzero
constant matrix in (20) is therefore not $\bar\partial$-exact.
No such $\psi$ exists. The smooth underlying bundle is fixed and trivial,
so its Chern classes vanish independently of this deformation problem.
Also $\operatorname{tr}[U,V]=0$: the obstruction is in a traceless
extension direction, not detected by its trace. This example rules out
the implication that Hodge Chern data and first-order compatibility solve
all the nonlinear integrability equations. It does not rule out a
different deformation direction or the trivial holomorphic bundle itself.
\end{proof}

For an actual source bundle the analogous equation is still
$\bar\partial\Theta+\Theta\wedge\Theta=0$ in its endomorphism-valued
Dolbeault complex. A fold's charge/period equation is a constraint on
cohomology; it supplies neither the missing traceless obstruction
cancellation nor the algebraic descent of a merely Stein holomorphic
bundle. On a projective realization, GAGA applies after a genuine
holomorphic coherent source has been constructed. It does not solve
equation (20). Thus the integrable-source and source-deformation proposals
must provide additional input at this nonlinear step.

\section{Typed GST chain and external leaves}
The chain actually used is:
\[
 (X,p,\alpha)\longmapsto(J,K,a,F_a)
 \xrightarrow[\text{proved only when supplied}]{\text{matched factor or }c_A\ne0}
 [\mathcal M_g]\text{ or }W_A
 \longmapsto Z_\alpha\xrightarrow{\operatorname{cl}^{\rm can}}\alpha.
\]
The first arrow is the previously derived topological/Stein construction.
The middle objects in (3) and (12) are actual algebraic sources; the last
two arrows are the genuine $K_0$/Chow normalization and affine-bundle
original-X return, with canonical Chern comparison. They do not use a
Stage-2G map chosen independently of canonical Betti semantics. No numerical
Gram identity is used as a Chow equality. Native Gram closure is unnecessary
when (4) or (13) applies; for arbitrary further sources it remains the
previous conditional finite-source operation, not their existence proof.

Authenticated external leaves checked on 3 November 2026 in the current
Stacks Project pages are:
\begin{itemize}
\item \href{https://stacks.math.columbia.edu/tag/00OE}{10.110.8, tag 00OE}
and \href{https://stacks.math.columbia.edu/tag/0656}{15.76.3, tag 0656}:
regular finite-dimensional affine coordinate ring and finite module imply
a finite projective resolution. These apply to (3), not to a singular base.
\item \href{https://stacks.math.columbia.edu/tag/037B}{10.37.1, tag 037B}:
normal domains are integrally closed. The factor-coefficient argument above
is provided here, rather than attributed to an unstated descent theorem.
\item \href{https://stacks.math.columbia.edu/tag/02UK}{42.43.1, tag 02UK}:
the line-filtration pullback is injective on Chow groups. This is the
splitting input in (10).
\item \href{https://stacks.math.columbia.edu/tag/0FA8}{42.44.1--2, tag 0FA8}
and \href{https://stacks.math.columbia.edu/tag/0FBI}{42.54.3--5, tag 0FBI}:
regular zero-section and excess-normal-bundle formulas. The explicit
normal inclusion above verifies their hypotheses for (6) and (16).
\end{itemize}
For (22)--(23), Pierre Deligne,
\href{https://www.claymath.org/wp-content/uploads/2022/06/hodge.pdf}{\emph{The Hodge Conjecture}},
official Clay problem description, Section 1, p.1, gives the compact
Kahler harmonic-type decomposition and filtration used here. This exact
page was opened and checked. The Cech contraction, its filtration behavior,
the cochain lift and telescoping computation are derived above.
The earlier project records remain the leaves for Jouanolou, topological
rational Chern character, Stein realization, and canonical comparison.
They are retained with their hand-proof status. No claim that all these
geometric structures were instantiated in Lean is made. These classical
inputs are not attributed as new OpenAI discoveries. No novelty search
or global novelty claim is part of this work.

\section{Focused adversarial audit and exact outstanding obligation}
The factor-to-cycle theorem is proved with its stated matched-factor
hypothesis; it permits resultant boundary and constructs the coherent
source there. The block/excess formulas include rank-zero selections,
$s=0,N$, positive-dimensional excess components, and all $p>0$.
The package construction requires the stated regular-incidence perturbation
and retains every local block. Root transitivity alone kills its package
coefficients; it does not prove vanishing of every finer boundary cycle.
The test (14) has a genuinely nonzero target class, not the zero-cluster
example of October 24. The source in (17) is conditional on algebraic
bundles, and auxiliary-hyperplane terms have not been discarded.
The obstruction (20) concerns a specified first-order direction and does
not disprove existence of all holomorphic sources.
The full filtered nerve comparison is proved for its explicitly constructed
realization; it is not claimed for every GST graph merely bearing the same
number of holes. Its cochain spaces can be infinite dimensional although
their Betti cohomology is finite dimensional. The rational structure is
retained through the constant nerve; the complex Hodge charge need not
preserve that rational lattice.

The universal conclusion remains incomplete at a geometric implication:
for the prescribed arbitrary higher-codimension Hodge input, no proof here
forces a matched factor, a nonzero computed package, a finer selected
cycle with the required class, or a solution of the source integrability
equations. The known Hodge premise of the provisional experiment cannot
be removed by any of the implications proved here. Consequently the
statement ``mathematics finished; only Lean remains'' is not established.
These are completed mathematical subproofs and executed route tests,
with a specific unproved universal leaf. Review provenance is self-review.
\end{document}
